\documentclass[UTF8]{ctexart}
\usepackage[a4paper,margin=2.4cm]{geometry}
\usepackage{booktabs}
\usepackage{longtable}
\usepackage{hyperref}
\usepackage{xcolor}
\usepackage{enumitem}

\hypersetup{
  colorlinks=true,
  linkcolor=blue,
  urlcolor=blue
}

\title{YOLO Final 双尺度三框实验诊断报告}
\author{yolo\_final 诊断记录}
\date{2026-04-28}

\begin{document}
\maketitle

\section{诊断范围}

本报告总结当前 \texttt{yolo\_final/} 中已经确认的问题、已完成修复、COCO-only 路线调整，以及后续训练实验建议。

本次诊断覆盖以下内容：

\begin{itemize}[leftmargin=2em]
  \item 双尺度三框正式实验 \texttt{dual\_scale\_three\_box\_formal\_20260427\_171851}；
  \item COCO 官方风格评估链路；
  \item \texttt{YOLOLoss} 中正样本统计、分类损失语义、多尺度 loss 聚合、动态匹配逻辑；
  \item 8GPU 训练没有被正确使用的问题，即表面多卡占用但实际计算没有高效分布；
  \item COCO-only 数据口径与后续实验配置；
  \item 当前已经启动的 \texttt{lambda\_noobj=1.0} 对照实验。
\end{itemize}

\section{已确认的关键问题}

\subsection{COCO 评估坐标空间错误}

最初的 COCO 评估结果异常接近 0：

\begin{center}
\begin{tabular}{lrr}
\toprule
指标 & 修复前 & 修复后 \\
\midrule
COCO AP & 0.000240 & 0.121699 \\
COCO AP50 & 0.000965 & 0.210424 \\
COCO AP75 & 0.000016 & 0.122554 \\
COCO AR@100 & 0.003765 & 0.263924 \\
\bottomrule
\end{tabular}
\end{center}

根因是预测框由 \texttt{decode\_predictions\_for\_image()} 解码为 resize 后的 $320\times320$ 坐标，而 COCO GT 使用原图坐标。旧版 \texttt{evaluate\_coco\_subset()} 直接把 resize 后坐标写入 pycocotools 结果，导致预测框与 GT 不在同一坐标系。

已修复位置：

\begin{itemize}[leftmargin=2em]
  \item \texttt{yolo\_final/utils/coco\_eval.py}
\end{itemize}

修复方式：

\begin{itemize}[leftmargin=2em]
  \item 从 dataset target 中读取 \texttt{original\_size}；
  \item 将预测框从 $320\times320$ 映射回原图宽高；
  \item 对坐标进行边界裁剪后再写入 COCO \texttt{bbox}。
\end{itemize}

验证：

\begin{itemize}[leftmargin=2em]
  \item 200 张 COCO val 的 GT-as-pred oracle 得到 AP/AP50 = 1.0；
  \item 修复后完整 COCO val 评估 AP50 从 0.000965 提升到 0.210424。
\end{itemize}

\subsection{VOC 与 COCO 混合标签空间冲突}

旧正式配置使用：

\begin{itemize}[leftmargin=2em]
  \item \texttt{all\_train.jsonl}
  \item \texttt{all\_val.jsonl}
  \item \texttt{num\_classes=80}
\end{itemize}

但 \texttt{all\_train/all\_val} 是 VOC 与 COCO 的直接拼接，VOC 使用 0--19 的 VOC-local 类别 id，COCO 使用 0--79 的 COCO-local 类别 id。同一个数字在两个数据集中语义不同。

\begin{center}
\begin{tabular}{rll}
\toprule
Label ID & VOC 语义 & COCO 语义 \\
\midrule
0 & aeroplane & person \\
2 & bird & car \\
6 & car & train \\
14 & person & bird \\
18 & train & sheep \\
19 & tvmonitor & cow \\
\bottomrule
\end{tabular}
\end{center}

影响：

\begin{itemize}[leftmargin=2em]
  \item 混合训练存在语义噪声；
  \item \texttt{all\_val} internal mAP 只能作为诊断指标，不适合作为正式主指标；
  \item COCO-only 评估是当前更可信的正式口径。
\end{itemize}

当前路线已经调整为 COCO-only：

\begin{itemize}[leftmargin=2em]
  \item \texttt{configs/dual\_scale\_three\_box\_coco\_only.toml}
  \item train: \texttt{coco2017\_train.jsonl}
  \item val: \texttt{coco2017\_val.jsonl}
\end{itemize}

\subsection{\texttt{positive\_cells\_per\_image} 统计维度错误}

旧代码：

\begin{verbatim}
positive_cells_per_image = object_mask.sum(dim=(1, 2)).mean()
\end{verbatim}

在多 anchor 情况下，\texttt{object\_mask} 形状是：

\begin{verbatim}
[B, H, W, A]
\end{verbatim}

旧写法在 sum 掉 H/W 后得到 \texttt{[B,A]}，再 mean 会把 anchor 维也平均掉。因此日志中的 \texttt{pos\_cells} 实际约为真实每图正样本数的 $1/A$。

已修复为：

\begin{verbatim}
positive_cells_per_image = object_mask.reshape(object_mask.shape[0], -1).sum(dim=1).mean()
\end{verbatim}

影响：

\begin{itemize}[leftmargin=2em]
  \item 不影响训练梯度；
  \item 影响 assignment 正样本数判断；
  \item 旧 run 的 \texttt{pos\_cells} 需要按 anchor 数量重新解释。
\end{itemize}

\subsection{\texttt{quality\_bce} 配置语义不显式}

配置中使用：

\begin{verbatim}
cls_loss_mode = "quality_bce"
soft_classification_target = "iou"
\end{verbatim}

旧代码只有 \texttt{varifocal} 分支，其余模式都走普通 BCE。因此 \texttt{quality\_bce} 实际含义是 ``IoU soft target + BCE''，但代码没有显式表达。

已修复：

\begin{itemize}[leftmargin=2em]
  \item \texttt{cls\_loss\_mode} 现在只允许 \texttt{bce}、\texttt{quality\_bce}、\texttt{varifocal}；
  \item \texttt{quality\_bce} 有显式分支；
  \item 未知模式会直接报错，避免配置拼写错误静默走 BCE。
\end{itemize}

\subsection{双尺度 loss 直接相加}

当前双尺度模型输出：

\begin{itemize}[leftmargin=2em]
  \item p4: $20\times20$
  \item p5: $10\times10$
\end{itemize}

旧逻辑对每个 scale 独立做 dynamic assignment，然后直接相加 loss。这会让同一个 GT 可能在 p4 和 p5 同时成为正样本，增加候选框密度。

本次没有默认改变行为，但新增了可配置能力：

\begin{itemize}[leftmargin=2em]
  \item \texttt{scale\_loss\_weights = "p4:1.0;p5:0.5"} 可对不同尺度 loss 加权；
  \item \texttt{scale\_assignment = "area"} 可按目标面积将 GT 分配给一个尺度；
  \item 默认仍为 \texttt{scale\_assignment="all"}，保持历史配置行为。
\end{itemize}

\subsection{\texttt{dynamic\_cost} 缺少 anchor shape prior}

旧动态匹配 cost：

\begin{verbatim}
cost = dynamic_box_cost * (1 - iou) + dynamic_cls_cost * cls_cost
\end{verbatim}

anchor 只参与预测框 decode，没有作为匹配先验。早期预测框较差时，IoU cost 可能不稳定。

本次新增可配置项：

\begin{verbatim}
dynamic_anchor_shape_cost = 0.0
\end{verbatim}

默认关闭，保持历史行为。若设置为正数，会加入轻量 shape ratio cost：

\begin{verbatim}
cost += dynamic_anchor_shape_cost * anchor_shape_cost
\end{verbatim}

\subsection{\texttt{dynamic\_topk=2} 在双尺度下可能偏多}

由于 p4 和 p5 独立匹配，每个 GT 在每个尺度最多可获得 \texttt{dynamic\_topk} 个正样本。当前配置 \texttt{dynamic\_topk=2} 时，一个 GT 理论上可在两个尺度共得到最多 4 个正样本。

这有利于召回，但也可能增加候选框数量和重复框。

本次新增候选配置：

\begin{itemize}[leftmargin=2em]
  \item \texttt{configs/dual\_scale\_three\_box\_coco\_only\_noobj1\_topk1.toml}
\end{itemize}

该配置不自动启动，仅作为后续对照实验入口。

\subsection{objectness 负样本权重偏低}

旧配置：

\begin{verbatim}
lambda_noobj = 0.5
\end{verbatim}

结合当前现象：

\begin{itemize}[leftmargin=2em]
  \item 候选框过多；
  \item 每图经常打满 \texttt{top\_k=100}；
  \item 可视化中红框过密；
  \item false positive 较多。
\end{itemize}

已经启动第一组训练方向对照：

\begin{itemize}[leftmargin=2em]
  \item 配置：\texttt{configs/dual\_scale\_three\_box\_coco\_only\_noobj1.toml}
  \item run id：\texttt{dual\_scale\_three\_box\_coco\_only\_noobj1\_20260428\_143454}
  \item 变量：\texttt{lambda\_noobj = 1.0}
\end{itemize}

注意：该训练进程在本次代码修复前已经启动，进程内使用的是启动时加载的旧 loss 代码。源码修复会影响后续新启动的实验，不会改变当前已运行进程。

\subsection{内部评估不是正式 COCO AP}

\texttt{utils/evaluation.py} 实现的是简化指标：

\begin{itemize}[leftmargin=2em]
  \item mAP@0.5
  \item precision
  \item recall
\end{itemize}

它适合快速诊断，但正式结论必须使用 \texttt{tools/evaluate\_coco.py} 的 pycocotools 指标：

\begin{itemize}[leftmargin=2em]
  \item AP@[0.50:0.95]
  \item AP50
  \item AP75
  \item APS/APM/APL
  \item AR@100
\end{itemize}

\subsection{\texttt{score\_alpha=2.0} 依赖 objectness 校准}

当前预测打分为：

\begin{verbatim}
score = obj_score ** score_alpha * cls_score ** score_beta
\end{verbatim}

配置中常用：

\begin{verbatim}
score_alpha = 2.0
score_beta = 1.0
\end{verbatim}

这会明显放大 objectness 校准质量对最终排序的影响。非训练诊断显示，提高阈值不能明显改善 AP，说明问题不只是简单阈值，而是排序质量和误检控制。

\section{本次代码修改汇总}

\begin{longtable}{lll}
\toprule
文件 & 修改内容 & 行为影响 \\
\midrule
\texttt{losses/yolo\_loss.py} & 修正 \texttt{positive\_cells\_per\_image} 统计 & 日志更准确 \\
\texttt{losses/yolo\_loss.py} & 显式支持 \texttt{quality\_bce} & 配置语义清晰 \\
\texttt{losses/yolo\_loss.py} & 新增 \texttt{dynamic\_anchor\_shape\_cost} & 默认关闭 \\
\texttt{losses/yolo\_loss.py} & 新增 \texttt{scale\_assignment} & 默认 all，不改变旧行为 \\
\texttt{losses/yolo\_loss.py} & 新增 \texttt{scale\_loss\_weights} & 默认每尺度权重 1.0 \\
\texttt{utils/config.py} & 配置摘要记录新字段 & 实验可追溯 \\
\texttt{tools/train.py} & 将新字段传入 \texttt{YOLOLoss} & 训练可配置 \\
\texttt{tools/smoke\_dual\_scale\_three\_box.py} & smoke 支持新字段 & 启动前可校验 \\
\texttt{engine/distributed\_trainer.py} & 新增 DDP 训练/验证循环 & 分布式指标汇总 \\
\texttt{tools/train\_ddp.py} & 新增 \texttt{torchrun} 入口 & 真正 8GPU 分布式训练 \\
\texttt{configs/...topk1.toml} & 新增 topk=1 候选配置 & 不自动启动 \\
\bottomrule
\end{longtable}

\section{COCO-only 当前状态}

COCO-only 数据已经准备：

\begin{center}
\begin{tabular}{lrrl}
\toprule
Split & Samples & Chunks & Packed Index \\
\midrule
train & 118287 & 116 & \texttt{coco2017\_train/chunk\_1024/index.json} \\
val & 5000 & 5 & \texttt{coco2017\_val/chunk\_1024/index.json} \\
\bottomrule
\end{tabular}
\end{center}

当前建议：

\begin{enumerate}[leftmargin=2em]
  \item 当前 \texttt{lambda\_noobj=1.0} run 可以继续完成，用来观察背景抑制方向；
  \item 因为该 run 启动早于本次 loss 代码修复，其 \texttt{pos\_cells} 日志仍按旧代码解释；
  \item 后续新实验应使用修复后的代码；
  \item 下一组候选实验优先考虑 \texttt{dynamic\_topk=1}，或开启轻量 \texttt{dynamic\_anchor\_shape\_cost}。
\end{enumerate}

\section{多卡训练未正确使用的诊断}

\subsection{问题定性}

该问题不是普通的“训练慢”，而是训练基础设施层面的多卡使用错误。

原训练脚本虽然能看到 8 张 GPU，也能让 8 张 GPU 分配显存，但实际训练并没有以正确的分布式方式运行。代码使用的是 \texttt{torch.nn.DataParallel}，而不是 \texttt{DistributedDataParallel}。在本项目当前的 heavy dynamic loss 设计下，这会导致：

\begin{itemize}[leftmargin=2em]
  \item 多张 GPU 主要参与模型 forward 的拆分；
  \item 模型输出会 gather 回主设备；
  \item \texttt{YOLOLoss}、dynamic assignment、指标统计等重计算没有真正按 GPU 分布式并行；
  \item GPU0 和单进程 Python 调度成为瓶颈；
  \item 其他 GPU 显存被占用，但计算利用率低。
\end{itemize}

因此，之前所谓“8GPU 训练”严格来说不是正确的 8GPU 分布式训练，而是低效的单进程 DataParallel 训练。这会影响实验效率，也会影响我们对训练成本、实验排期和可扩展性的判断。

\subsection{问题现象}

在 COCO-only + \texttt{lambda\_noobj=1.0} 实验中，虽然训练程序识别到了 8 张 GPU，且 8 张 GPU 都有显存占用，但单个 epoch 的训练时间明显偏长。日志中观测到：

\begin{center}
\begin{tabular}{lrrr}
\toprule
阶段 & Batch 数 & 用时 & 平均 step 时间 \\
\midrule
epoch 1 train & 925 & 2856.5s & 约 3.09s/step \\
epoch 2 train & 925 & 3034.7s & 约 3.28s/step \\
epoch 1 val & 40 & 131.3s & 约 3.28s/step \\
epoch 2 val & 40 & 125.0s & 约 3.13s/step \\
\bottomrule
\end{tabular}
\end{center}

当前训练配置为：

\begin{verbatim}
batch_size = 128
num_workers = 64
pin_memory = true
persistent_workers = true
prefetch_factor = 4
use_data_parallel = true
\end{verbatim}

按日志估算，训练吞吐大约只有 39--41 images/s。对于 8 张 RTX 4090 和 320 分辨率输入，这个速度偏慢。

\subsection{排查证据}

排查时观察到以下现象：

\begin{itemize}[leftmargin=2em]
  \item 训练没有卡死，step 日志持续推进；
  \item \texttt{.pt} 数据加载路径正常，训练集为 118287 张 COCO train，验证集为 5000 张 COCO val；
  \item DataLoader 配置已经较激进，\texttt{num\_workers=64}，且启用了 \texttt{pin\_memory}、\texttt{persistent\_workers} 和 \texttt{prefetch\_factor=4}；
  \item GPU 显存被 8 张卡占用，但 GPU 利用率不均衡，常见状态是 GPU0 较高，其他 GPU 很低；
  \item 进程结构显示当前训练实际是单个主 Python 进程配合大量 DataLoader worker，而不是 8 个独立训练进程。
\end{itemize}

因此，这次慢不是因为数据路径错误，也不是因为训练完全没有启动，而是多 GPU 训练方式效率低。

\subsection{根因分析}

当前 \texttt{tools/train.py} 中多 GPU 使用的是：

\begin{verbatim}
torch.nn.DataParallel(model)
\end{verbatim}

这和真正的 \texttt{DistributedDataParallel} 不同。DataParallel 是单进程多卡模式，主要问题包括：

\begin{itemize}[leftmargin=2em]
  \item GPU0 负责 scatter/gather，容易成为瓶颈；
  \item batch 被分发到各 GPU 做 forward，但输出会 gather 回主设备；
  \item 当前 loss 没有被分散到各 GPU 独立计算，而是在 gather 后集中计算；
  \item 动态匹配 loss 本身比较重，集中计算会进一步放大 GPU0 和 Python 调度瓶颈。
\end{itemize}

从 bug 角度看，错误点在于：配置字段 \texttt{use\_data\_parallel=true} 容易让人误以为已经在进行标准 8GPU 训练，但实际代码路径并没有启动 8 个训练进程，也没有使用 \texttt{DistributedSampler}，更没有让每张 GPU 独立完成 loss 计算。因此，本轮训练虽然可以产出结果，但不应作为“正确多卡训练流程”的模板继续复用。

本实验的 loss 尤其重，原因是：

\begin{itemize}[leftmargin=2em]
  \item 双尺度输出：\texttt{p4,p5}；
  \item 每个尺度 3 个 anchor；
  \item 使用 \texttt{assignment\_strategy = "dynamic\_cost"}；
  \item 每张图、每个 GT 都需要做候选筛选、IoU/cost 计算和 top-k 分配；
  \item \texttt{dynamic\_topk = 2} 时，双尺度下正样本候选数量进一步增加。
\end{itemize}

因此，实际瓶颈是：

\begin{quote}
模型 forward 看似用了 8 张 GPU，但 heavy loss / dynamic assignment 没有真正按 GPU 并行拆开，导致 8GPU 利用率很差。
\end{quote}

\subsection{解决方案}

为解决该问题，新增了 DDP 训练入口和分布式训练循环：

\begin{itemize}[leftmargin=2em]
  \item \texttt{tools/train\_ddp.py}
  \item \texttt{engine/distributed\_trainer.py}
\end{itemize}

该修复不是为了改变模型结构或 loss 公式，而是修复训练执行方式：后续正式多卡实验必须通过 \texttt{torchrun} 启动 DDP 脚本，而不能继续把 \texttt{tools/train.py} 的 DataParallel 路径当作 8GPU 正式训练路径。

新的 DDP 方案使用 \texttt{torchrun} 启动，每张 GPU 一个进程。每个进程独立完成：

\begin{itemize}[leftmargin=2em]
  \item 本地 batch 读取；
  \item 本地 model forward；
  \item 本地 \texttt{YOLOLoss} 和 dynamic assignment；
  \item 本地 backward；
  \item DDP 梯度同步。
\end{itemize}

这样 dynamic assignment 和 loss 计算会被分摊到 8 个进程/8 张 GPU，而不是集中在 DataParallel 的主设备上。

\subsection{DDP 脚本设计}

新增脚本的关键设计如下：

\begin{itemize}[leftmargin=2em]
  \item 使用 \texttt{torch.distributed.init\_process\_group} 初始化；
  \item CUDA 环境下默认使用 \texttt{nccl} backend；
  \item 使用 \texttt{DistributedDataParallel} 包装模型；
  \item 使用 \texttt{DistributedSampler} 分割训练集和验证集；
  \item 只有 rank0 写 TensorBoard、保存 checkpoint、更新 metadata、写 result summary 和生成可视化；
  \item epoch 指标通过 \texttt{dist.all\_reduce} 汇总；
  \item 默认将配置里的 \texttt{train.batch\_size} 解释为全局 batch size；
  \item 默认将配置里的 \texttt{train.num\_workers} 按 world size 拆分。
\end{itemize}

例如当前配置：

\begin{verbatim}
batch_size = 128
num_workers = 64
\end{verbatim}

在 8GPU DDP 下会被解释为：

\begin{center}
\begin{tabular}{lrr}
\toprule
项目 & 配置值 & DDP 实际值 \\
\midrule
global batch size & 128 & 128 \\
per-rank batch size & - & 16 \\
global workers & 64 & 64 \\
per-rank workers & - & 8 \\
\bottomrule
\end{tabular}
\end{center}

这样可以避免 \texttt{torchrun} 后每个 rank 都启动 64 个 worker，导致总 worker 数膨胀到 512。

\subsection{已完成验证}

新增脚本已经完成以下验证：

\begin{enumerate}[leftmargin=2em]
  \item \texttt{python -m py\_compile engine/distributed\_trainer.py tools/train\_ddp.py} 通过；
  \item \texttt{python tools/train\_ddp.py --help} 通过；
  \item CPU 单进程 1-step smoke run 通过；
  \item CPU \texttt{torchrun --nproc\_per\_node=2} 的真实 DDP smoke run 通过；
  \item 两个 smoke run 均使用 COCO-only \texttt{.pt} 数据，并完成 1 个 train step 和 1 个 val step；
  \item smoke run 正常写出 run directory、metadata、checkpoint 和 result summary。
\end{enumerate}

验证用的小测试没有占用当前正在运行的 8GPU DataParallel 训练。

\subsection{正式 DDP 启动命令}

后续若要启动正式 8GPU DDP 训练，使用：

\begin{verbatim}
cd /home/lidz/YOLO/yolo_final
tmux new -d -s yolo_coco_noobj1_ddp_8gpu \
'source /home/lidz/miniconda3/etc/profile.d/conda.sh && \
conda activate yolov1 && \
CUDA_VISIBLE_DEVICES=0,1,2,3,4,5,6,7 \
torchrun --standalone --nproc_per_node=8 \
tools/train_ddp.py \
--config configs/dual_scale_three_box_coco_only_noobj1.toml \
2>&1 | tee logs/dual_scale_three_box_coco_only_noobj1_ddp_8gpu_tmux.log'
\end{verbatim}

监控命令：

\begin{verbatim}
tmux attach -t yolo_coco_noobj1_ddp_8gpu
tail -f logs/dual_scale_three_box_coco_only_noobj1_ddp_8gpu_tmux.log
\end{verbatim}

\subsection{注意事项}

当前 DataParallel 训练仍在运行，尚未被 DDP 替换。若要正式切换到 DDP，应先停止当前 DataParallel run，记录其状态，然后使用上述命令启动新的 DDP run。两个 run 的 run name 会不同：

\begin{itemize}[leftmargin=2em]
  \item 当前 DataParallel：\texttt{dual\_scale\_three\_box\_coco\_only\_noobj1\_*}
  \item 新 DDP：\texttt{dual\_scale\_three\_box\_coco\_only\_noobj1\_ddp\_*}
\end{itemize}

因此两者不会覆盖输出目录、TensorBoard 目录或记录目录。

\section{COCO-only + \texttt{lambda\_noobj=1.0} 完成结果}

\subsection{训练状态}

本轮固定代码后的 DataParallel 训练已经完成：

\begin{itemize}[leftmargin=2em]
  \item Run ID：\texttt{dual\_scale\_three\_box\_coco\_only\_noobj1\_20260428\_180817}
  \item 配置：\texttt{configs/dual\_scale\_three\_box\_coco\_only\_noobj1.toml}
  \item 数据：COCO-only \texttt{.pt}
  \item 状态：\texttt{completed}
  \item last epoch：20
  \item best epoch：20
  \item best val loss：5.262633669376373
  \item global step：18500
\end{itemize}

训练和验证 loss 持续下降，关键 epoch 如下：

\begin{center}
\begin{tabular}{rrrrrrrr}
\toprule
Epoch & Train Total & Train Box & Train Obj & Train Cls & Val Total & Val Box & Val Obj \\
\midrule
1 & 7.498224 & 1.104558 & 1.904885 & 0.070552 & 7.006711 & 1.009949 & 1.901616 \\
10 & 5.448432 & 0.751085 & 1.648955 & 0.044049 & 5.733664 & 0.796673 & 1.704529 \\
20 & 4.707622 & 0.633448 & 1.500294 & 0.040087 & 5.262634 & 0.723287 & 1.604132 \\
\bottomrule
\end{tabular}
\end{center}

该 run 的 train loop 总耗时约 18.37 小时，validation loop 总耗时约 0.73 小时，合计约 19.11 小时。这进一步验证了上一节关于 DataParallel 训练效率不足的诊断。

\subsection{正式 COCO 指标}

使用 \texttt{best.pth} 在完整 COCO val2017 5000 张图上进行 pycocotools 评估，结果如下：

\begin{center}
\begin{tabular}{lr}
\toprule
Metric & Value \\
\midrule
AP@[.50:.95] & 0.126985 \\
AP50 & 0.218369 \\
AP75 & 0.129184 \\
AR1 & 0.166430 \\
AR10 & 0.251601 \\
AR100 & 0.260567 \\
Predictions & 206786 \\
Predictions/Image & 41.36 \\
\bottomrule
\end{tabular}
\end{center}

与上一轮 corrected formal run 对照：

\begin{center}
\begin{tabular}{lrrrr}
\toprule
Run & AP & AP50 & AP75 & AR100 \\
\midrule
previous corrected formal & 0.121699 & 0.210424 & 0.122554 & 0.263924 \\
COCO-only + \texttt{lambda\_noobj=1.0} & 0.126985 & 0.218369 & 0.129184 & 0.260567 \\
delta & +0.005286 & +0.007945 & +0.006630 & -0.003357 \\
\bottomrule
\end{tabular}
\end{center}

结论是：\texttt{lambda\_noobj=1.0} 带来小幅 AP/AP50/AP75 提升，但 AR100 略降。这符合更强背景抑制的预期：误检排序略有改善，但召回略受影响。由于上一轮 formal run 仍受 mixed VOC+COCO 标签空间问题影响，该对照只能作为方向性参考，而不是严格消融。

\subsection{诊断扫描}

在 500 张 COCO val 子集上做了非训练诊断扫描。默认设置为：

\begin{verbatim}
score_threshold = 0.05
score_alpha = 2.0
score_beta = 1.0
top_k = 100
nms_iou_threshold = 0.5
\end{verbatim}

诊断显示候选仍然偏多：

\begin{itemize}[leftmargin=2em]
  \item diagnostic prediction count：50000
  \item predictions per image：min 100，mean 100，max 100
  \item 说明候选列表仍然打满 \texttt{top\_k=100}
\end{itemize}

score threshold 扫描：

\begin{center}
\begin{tabular}{rrrrr}
\toprule
Threshold & AP & AP50 & AR100 & Predictions \\
\midrule
0.01 & 0.157420 & 0.254476 & 0.300076 & 35660 \\
0.03 & 0.157235 & 0.253744 & 0.296141 & 26643 \\
0.05 & 0.156741 & 0.252505 & 0.287816 & 20875 \\
0.10 & 0.154733 & 0.247743 & 0.265047 & 12416 \\
0.20 & 0.150149 & 0.237212 & 0.234355 & 5919 \\
0.30 & 0.135472 & 0.210205 & 0.183123 & 3311 \\
\bottomrule
\end{tabular}
\end{center}

top-k 扫描：

\begin{center}
\begin{tabular}{rrrrr}
\toprule
Top-k & AP & AP50 & AR100 & Predictions \\
\midrule
20 & 0.153218 & 0.242610 & 0.250194 & 9000 \\
50 & 0.155875 & 0.250427 & 0.276953 & 16618 \\
100 & 0.156741 & 0.252505 & 0.287816 & 20875 \\
200 & 0.156814 & 0.252705 & 0.289030 & 22009 \\
\bottomrule
\end{tabular}
\end{center}

这说明 post-processing 只能带来很小变化。降低阈值会略增 AP，但预测数量明显增加；提高阈值会减少预测，但 AP/AR 同时下降。下一步仍应从训练侧改善 assignment 和 ranking，而不是只调推理阈值。

\subsection{本轮产物}

\begin{itemize}[leftmargin=2em]
  \item Summary：\texttt{logs/dual\_scale\_three\_box\_coco\_only\_noobj1\_20260428\_180817\_summary.md}
  \item COCO eval：\texttt{outputs/evaluations/dual\_scale\_three\_box\_coco\_only\_noobj1\_20260428\_180817\_coco\_eval.json}
  \item Diagnostics：\texttt{outputs/evaluations/dual\_scale\_three\_box\_coco\_only\_noobj1\_20260428\_180817\_diagnostics.json}
  \item Diagnostic visualizations：\texttt{outputs/evaluations/dual\_scale\_three\_box\_coco\_only\_noobj1\_20260428\_180817\_diagnostic\_vis/}
\end{itemize}

\section{COCO-only + \texttt{lambda\_noobj=1.0} + \texttt{dynamic\_topk=1} DDP 完成结果}

\subsection{实验目的}

本轮实验用于验证：将动态匹配从 \texttt{dynamic\_topk=2} 收紧到 \texttt{dynamic\_topk=1} 后，是否能够减少双尺度三框训练中的正样本压力、重复候选和误检密度。

同时，本轮也是修复多卡训练路径后的第一轮正式 8GPU DDP 完整训练，用来确认 \texttt{tools/train\_ddp.py} 和 \texttt{engine/distributed\_trainer.py} 能作为后续正式实验基础。

\subsection{训练状态}

本轮 DDP 训练已经完成：

\begin{itemize}[leftmargin=2em]
  \item Run ID：\texttt{dual\_scale\_three\_box\_coco\_only\_noobj1\_topk1\_ddp\_20260429\_144209}
  \item 配置：\texttt{configs/dual\_scale\_three\_box\_coco\_only\_noobj1\_topk1.toml}
  \item 数据：COCO-only \texttt{.pt}
  \item 启动方式：\texttt{torchrun} + DDP
  \item world size：8
  \item 每卡 batch size：16
  \item 全局 batch size：128
  \item best epoch：20
  \item best val loss：5.148357025909424
\end{itemize}

最终 epoch 指标如下：

\begin{center}
\begin{tabular}{lrrrrrrrr}
\toprule
Split & Total & Box & Obj & Cls & GIoU & Obj Target & Pos Cells & Dropped GT \\
\midrule
train e20 & 4.765897 & 0.627729 & 1.584755 & 0.042496 & 0.686135 & 0.700054 & 13.993169 & 0.547832 \\
val e20 & 5.148357 & 0.693390 & 1.637210 & 0.044196 & 0.653305 & 0.671206 & 14.097600 & 0.614800 \\
\bottomrule
\end{tabular}
\end{center}

本轮 train loop 总耗时 16008.707 秒，validation loop 总耗时 329.373 秒，合计约 4.538 小时。相比上一轮 DataParallel 约 19.11 小时，DDP 训练路径显著加快，说明此前“8GPU 没有正确高效训练”的诊断是成立的。

\subsection{正式 COCO 指标}

使用 \texttt{best.pth} 在完整 COCO val2017 5000 张图上进行 pycocotools 评估：

\begin{center}
\begin{tabular}{lr}
\toprule
Metric & Value \\
\midrule
AP@[.50:.95] & 0.108180 \\
AP50 & 0.195980 \\
AP75 & 0.105638 \\
AR1 & 0.151326 \\
AR10 & 0.229976 \\
AR100 & 0.239071 \\
Predictions & 244795 \\
Predictions/Image & 48.96 \\
\bottomrule
\end{tabular}
\end{center}

与上一轮 COCO-only + \texttt{lambda\_noobj=1.0} + \texttt{dynamic\_topk=2} 对照：

\begin{center}
\begin{tabular}{lrrr}
\toprule
Metric & topk=2 & topk=1 DDP & Delta \\
\midrule
Best val loss & 5.262634 & 5.148357 & -0.114277 \\
AP & 0.126985 & 0.108180 & -0.018805 \\
AP50 & 0.218369 & 0.195980 & -0.022389 \\
AP75 & 0.129184 & 0.105638 & -0.023546 \\
AR100 & 0.260567 & 0.239071 & -0.021496 \\
Val pos cells & 27.332 & 14.098 & -13.234 \\
\bottomrule
\end{tabular}
\end{center}

结论是：\texttt{dynamic\_topk=1} 的确减少了正样本数量，并使 val loss 更低；但正式 COCO AP、AP50、AP75、AR100 全部下降。这说明更严格的 assignment 并没有改善检测质量，上一轮 \texttt{topk=2} 对召回和排序有实际帮助。因此 \texttt{dynamic\_topk=1} 不应作为下一轮质量改进的默认方向。

\subsection{诊断扫描}

在 500 张 COCO val 子集上做了诊断扫描。默认推理设置仍为：

\begin{verbatim}
score_threshold = 0.05
score_alpha = 2.0
score_beta = 1.0
top_k = 100
nms_iou_threshold = 0.5
\end{verbatim}

raw diagnostic prediction 统计：

\begin{center}
\begin{tabular}{lr}
\toprule
Statistic & Value \\
\midrule
Prediction count & 50000 \\
Predictions/Image mean & 100.0 \\
Predictions/Image min & 100 \\
Predictions/Image max & 100 \\
Score q50 & 0.048543 \\
Score q75 & 0.117297 \\
Score q90 & 0.244761 \\
Score q95 & 0.371668 \\
Score q99 & 0.687747 \\
\bottomrule
\end{tabular}
\end{center}

top-k 扫描：

\begin{center}
\begin{tabular}{rrrrr}
\toprule
Top-k & AP & AP50 & AR100 & Predictions \\
\midrule
20 & 0.135333 & 0.224608 & 0.225878 & 9229 \\
50 & 0.138710 & 0.233835 & 0.252944 & 18153 \\
100 & 0.139413 & 0.235948 & 0.263971 & 24605 \\
200 & 0.139516 & 0.236225 & 0.266481 & 27220 \\
\bottomrule
\end{tabular}
\end{center}

score threshold 扫描：

\begin{center}
\begin{tabular}{rrrrr}
\toprule
Threshold & AP & AP50 & AR100 & Predictions \\
\midrule
0.01 & 0.139998 & 0.237598 & 0.275860 & 39642 \\
0.03 & 0.139759 & 0.236844 & 0.270529 & 30905 \\
0.05 & 0.139413 & 0.235948 & 0.263971 & 24605 \\
0.10 & 0.137979 & 0.232019 & 0.245194 & 14698 \\
0.20 & 0.132834 & 0.219203 & 0.211559 & 6616 \\
0.30 & 0.124544 & 0.201703 & 0.173921 & 3656 \\
\bottomrule
\end{tabular}
\end{center}

诊断结论：

\begin{itemize}[leftmargin=2em]
  \item \texttt{dynamic\_topk=1} 并没有解决 raw 候选打满 \texttt{top\_k=100} 的问题；
  \item 推理 \texttt{top\_k} 从 100 增加到 200，AP 几乎不变，说明当前不是简单截断不足；
  \item 提高 score threshold 会减少预测数量，但 AP 和 AR 同时下降；
  \item 下一步应回到 \texttt{dynamic\_topk=2}，优先尝试 anchor shape prior、objectness/ranking 校准或尺度分配，而不是继续收紧 top-k。
\end{itemize}

\subsection{本轮产物}

\begin{itemize}[leftmargin=2em]
  \item Summary：\texttt{logs/dual\_scale\_three\_box\_coco\_only\_noobj1\_topk1\_ddp\_20260429\_144209\_summary.md}
  \item COCO eval：\texttt{outputs/evaluations/dual\_scale\_three\_box\_coco\_only\_noobj1\_topk1\_ddp\_20260429\_144209\_coco\_eval.json}
  \item Diagnostics：\texttt{outputs/evaluations/dual\_scale\_three\_box\_coco\_only\_noobj1\_topk1\_ddp\_20260429\_144209\_diagnostics.json}
  \item Diagnostic visualizations：\texttt{outputs/evaluations/dual\_scale\_three\_box\_coco\_only\_noobj1\_topk1\_ddp\_20260429\_144209\_diagnostic\_vis/}
\end{itemize}

\section{后续实验建议}

\subsection{最小对照顺序}

\begin{enumerate}[leftmargin=2em]
  \item 保留 DDP 作为后续正式训练路径；
  \item 不把 \texttt{dynamic\_topk=1} 作为默认质量改进方案；
  \item 回到 \texttt{dynamic\_topk=2}；
  \item 优先尝试 \texttt{dynamic\_anchor\_shape\_cost=0.25/0.5}；
  \item 再考虑 objectness/ranking 校准，例如重新扫描 \texttt{score\_alpha} 或调整 objectness target；
  \item 最后再考虑按面积分配尺度或 scale loss weights。
\end{enumerate}

\subsection{正式评估标准}

每轮正式实验完成后必须记录：

\begin{itemize}[leftmargin=2em]
  \item COCO AP/AP50/AP75；
  \item AR@100；
  \item prediction count；
  \item score quantiles；
  \item top predicted classes；
  \item 固定 COCO val 样例可视化；
  \item 配置快照和 run id。
\end{itemize}

\section{结论}

当前最大的问题不是模型无法训练，而是：

\begin{itemize}[leftmargin=2em]
  \item 旧 COCO 评估坐标错误已经修复；
  \item 旧混合数据标签空间不干净，正式路线已转为 COCO-only；
  \item 旧 \texttt{pos\_cells} 日志低估正样本数；
  \item \texttt{quality\_bce} 语义已经显式化；
  \item 8GPU 训练慢的主要原因是 DataParallel 无法高效并行 heavy dynamic loss，已补充 DDP 训练脚本；
  \item 双尺度重复监督、动态 top-k 和缺少 anchor prior 是后续主要实验变量。
\end{itemize}

因此，后续判断实验成败应以 COCO-only 的正式 COCO 指标和固定诊断可视化为准，而不是仅以训练 loss 或旧 internal mAP 为准。

\end{document}
